\documentclass[a4paper,11pt]{article}
\usepackage[utf8]{inputenc}
\usepackage[T1]{fontenc}
\usepackage[french]{babel}
\usepackage[top=2cm, bottom=2cm, left=2.2cm, right=2.2cm]{geometry}
\usepackage{xcolor}
\usepackage{graphicx}
\usepackage{booktabs}
\usepackage{tabularx}
\usepackage{enumitem}
\usepackage{listings}
\usepackage{float}
\usepackage{caption}
\usepackage{amssymb}
\usepackage{pifont}
\usepackage{tcolorbox}
\usepackage{hyperref}

\hypersetup{colorlinks=true, linkcolor=primary, urlcolor=primary}

\definecolor{primary}{HTML}{1A237E}
\definecolor{success}{HTML}{2E7D32}
\definecolor{warning}{HTML}{E65100}
\definecolor{danger}{HTML}{B71C1C}
\definecolor{lightgreen}{HTML}{E8F5E9}
\definecolor{lightorange}{HTML}{FFF3E0}
\definecolor{lightblue}{HTML}{E3F2FD}
\definecolor{codebg}{HTML}{F5F5F5}
\definecolor{codeblue}{rgb}{0.13,0.29,0.53}
\definecolor{codegreen}{rgb}{0.0,0.50,0.0}
\definecolor{codegray}{rgb}{0.5,0.5,0.5}

\lstdefinestyle{sh}{
  backgroundcolor=\color{codebg},
  basicstyle=\ttfamily\footnotesize,
  keywordstyle=\color{codeblue}\bfseries,
  commentstyle=\color{codegreen}\itshape,
  numberstyle=\tiny\color{codegray},
  breaklines=true,
  frame=single,
  showstringspaces=false,
  tabsize=2,
  columns=fullflexible,
}
\lstset{style=sh}

\newcommand{\ok}{\textcolor{success}{\checkmark}}
\newcommand{\ko}{\textcolor{danger}{\ding{55}}}

% ── Images à placer dans le même dossier que ce .tex ──────────────────
%   img_protova1.jpg → Protova1 (châssis jaune, Pico visible)
%   img_protova2.jpg → Protova2 (châssis rouge, Jetson + caméra + antennes 5G)
%   img_g29.jpg      → Volant G29 + pédalier sur support
%   img_ps4.jpg      → Manette PS4 blanche
%   img_rqt.png      → Capture rqt (flux caméra + liste topics)
% ──────────────────────────────────────────────────────────────────────

\title{\textbf{Rapport technique hebdomadaire --- ProtoVA2}\\[2mm]
\large Démonstration WiFi, mise en service du LiDAR\\ et conduite autonome réactive}
\author{Douae Sebti}
\date{Semaine du 16 au 19 juin 2026}

\begin{document}
\maketitle

\begin{abstract}
\noindent Ce rapport regroupe les travaux de la semaine sur le véhicule ProtoVA2
(Jetson Orin Nano, ROS\,2 Humble), en vue de la conduite autonome. Quatre volets
sont traités : (1) une démonstration de \textbf{téléopération sur WiFi} en DDS
natif (sans Zenoh), accompagnée d'outils de \textbf{supervision et de benchmark}
temps réel du lien réseau ; (2) la \textbf{conception d'un module de conduite autonome
réactive} de type \emph{follow-the-gap}, intégré à l'arbitrage \texttt{twist\_mux}
et validé hors matériel par un scan synthétique ; (3) la \textbf{mise en service
du LiDAR}, présentée comme un journal de diagnostic incluant volontairement les
tentatives infructueuses, jusqu'à la découverte et la réparation d'une panne
matérielle, puis la visualisation du flux \texttt{/scan} dans RViz ; (4) un
\textbf{premier essai de conduite réactive} (roues en l'air) validant la chaîne
de bout en bout. L'environnement, les conventions, les scripts livrés et les
prochaines étapes sont détaillés.
\end{abstract}

\tableofcontents
\bigskip

% =====================================================================
\section{Contexte et objectifs}
Le ProtoVA2 fonctionne sous ROS\,2 Humble (\texttt{ROS\_DOMAIN\_ID=125}). La
téléopération 5G existante repose sur \texttt{rmw\_zenoh\_cpp} (Zenoh en
hub-and-spoke), car le réseau cellulaire bloque le multicast UDP. Les objectifs
de la semaine étaient : préparer une \textbf{démonstration sur WiFi}, puis
\textbf{démarrer le volet autonomie} (perception LiDAR et conduite réactive).
L'essentiel du travail a été réalisé à distance, par SSH depuis le PC vers la
Jetson (\texttt{protova2@10.37.11.28}).

% =====================================================================
\section{Plateformes matérielles}
Les travaux portent sur deux prototypes de voiture RC autonome développés à
ALTEN Labs. \textbf{Protova1} est une reconstruction depuis zéro (phase
d'intégration) ; \textbf{Protova2} est la plateforme complète utilisée pour la
téléopération et l'autonomie (Jetson, caméra Orbbec, antennes 5G). Les deux
utilisent un Raspberry Pi Pico (RP2040) pour la commande PWM des actionneurs.

\begin{figure}[H]
\centering
\begin{minipage}[b]{0.47\textwidth}
  \includegraphics[width=\linewidth]{img_protova1.jpg}
  \caption{Protova1 --- reconstruction depuis zéro (châssis jaune, Raspberry Pi Pico RP2040 visible, phase d'intégration).}
\end{minipage}
\hfill
\begin{minipage}[b]{0.47\textwidth}
  \includegraphics[width=\linewidth]{img_protova2.jpg}
  \caption{Protova2 --- plateforme complète (châssis rouge, Jetson Nano, caméra Orbbec, antennes 5G, Seeed Studio).}
\end{minipage}
\end{figure}

\noindent La téléopération s'appuie sur plusieurs sources de commande, arbitrées
par \texttt{twist\_mux} (cf.\ §\ref{sec:mux}) : le volant Logitech \textbf{G29}
(USB, avec retour de force) et la manette \textbf{PS4} (Bluetooth, alternative
sans fil).

\begin{figure}[H]
\centering
\begin{minipage}[b]{0.47\textwidth}
  \includegraphics[width=\linewidth]{img_g29.jpg}
  \caption{Volant G29 + pédalier --- poste de pilotage ALTEN Labs.}
\end{minipage}
\hfill
\begin{minipage}[b]{0.47\textwidth}
  \includegraphics[width=\linewidth]{img_ps4.jpg}
  \caption{Manette PS4 --- téléopération Bluetooth (source secondaire, priorité 60 dans \texttt{twist\_mux}).}
\end{minipage}
\end{figure}

% =====================================================================
\section{Démonstration WiFi (DDS natif)}
Sur un réseau local WiFi, la découverte ROS\,2 par multicast fonctionne : Zenoh
n'est donc pas nécessaire. On utilise le middleware par défaut
\texttt{rmw\_fastrtps\_cpp}.
\begin{itemize}[leftmargin=1.4em]
  \item PC : \texttt{10.37.11.32} \quad--\quad Jetson : \texttt{10.37.11.28}
        (même sous-réseau \texttt{10.37.11.0/24}).
  \item Variables communes à tous les scripts WiFi :
\begin{lstlisting}[language=bash]
export ROS_DOMAIN_ID=125
export RMW_IMPLEMENTATION=rmw_fastrtps_cpp
unset ZENOH_SESSION_CONFIG_URI ZENOH_ROUTER_CONFIG_URI ZENOH_CONFIG
\end{lstlisting}
\end{itemize}

\subsection{Scripts créés}
\begin{table}[H]
\centering
\begin{tabularx}{\textwidth}{@{}l l X@{}}
\toprule
\textbf{Machine} & \textbf{Fichier} & \textbf{Rôle} \\
\midrule
PC     & \texttt{\textasciitilde/protova2/launch\_pc\_wifi.sh}  & \texttt{joy\_node} + retour de force G29 \\
Jetson & \texttt{\textasciitilde/launch\_car\_wifi.sh}          & \texttt{g29\_teleop}, TX, RX, \texttt{twist\_mux}, caméra Orbbec \\
PC     & \texttt{\textasciitilde/protova2/view\_camera\_wifi.sh}& Visualiseur caméra (\texttt{rqt\_image\_view}) \\
\bottomrule
\end{tabularx}
\end{table}
\noindent Tous reprennent les conventions des scripts 5G (nettoyage des nodes
résiduels, \texttt{trap}+\texttt{wait} pour un arrêt propre au \texttt{Ctrl+C}),
mais \textbf{sans Zenoh ni configuration 5G} (pas de dongle, pas de correctif
d'adressage /30$\rightarrow$/28, pas de watchdog).

\subsection{Validation}
Test de bout en bout : le PC découvre l'ensemble des topics et nodes du véhicule
via FastDDS, et les données circulent réellement.
\begin{center}
\begin{tabular}{@{}ll@{}}
\toprule
Vérification & Résultat \\
\midrule
Découverte des 35 topics / 7 nodes & \textcolor{success}{OK} \\
\texttt{/imu} (Jetson $\rightarrow$ PC) & $\sim$37--59 Hz \\
\texttt{/camera/color/camera\_info} & 15 Hz (= 15 fps configurés) \\
\texttt{/velocity} & message reçu \\
Isolation du point d'accès (AP isolation) & aucune \\
\bottomrule
\end{tabular}
\end{center}

\subsection{Correction d'un plantage de \texttt{rqt} sur la caméra}
En sélectionnant directement le topic \texttt{/camera/color/image\_raw/compressed},
\texttt{rqt} plante :
\begin{lstlisting}
create_subscription() ... existing topic ... incompatible type CompressedImage_
-> invalid allocator -> terminate (RCLError)
\end{lstlisting}
\textbf{Cause :} les topics \texttt{.../compressed}, \texttt{.../theora}, etc.\ ne
sont pas des topics indépendants mais des \emph{transports} du topic de base. Il
faut s'abonner au topic \textbf{de base} et choisir le transport :
\begin{lstlisting}[language=bash]
ros2 run rqt_image_view rqt_image_view /camera/color/image_raw \
  _image_transport:=compressed
\end{lstlisting}
C'est ce que fait le script \texttt{view\_camera\_wifi.sh} (le transport
\texttt{compressed} est en outre recommandé sur WiFi, plus léger que le brut
1280$\times$720).

% =====================================================================
\section{Supervision et benchmark du lien réseau}
Pour qualifier la qualité du lien (latence, débit, pertes), deux tableaux de bord
web temps réel ont été développés cette semaine (Python, bibliothèque standard
uniquement), servis par la Jetson et consultés depuis le navigateur du PC.
\begin{table}[H]
\centering
\begin{tabularx}{\textwidth}{@{}l X l@{}}
\toprule
\textbf{Outil} & \textbf{Rôle} & \textbf{Port} \\
\midrule
\texttt{monitor\_5g.py}   & \emph{monitoring passif} : latence (ICMP + RTT TCP Zenoh), débit RX/TX (Mbps), charge réseau (paquets/s) & 8088 \\
\texttt{bench\_5g\_live.py} & \emph{benchmark actif} : tests \texttt{iperf3} périodiques UL+DL, latence ICMP, débit TCP, perte paquets, bouton « Tester maintenant » & 8089 \\
\bottomrule
\end{tabularx}
\end{table}
\noindent Le benchmark actif (\texttt{bench\_5g\_live.py}) requiert un
\texttt{iperf3 -s} côté PC. Référence mesurée : $\approx$\,635 Mbps en descente /
70 Mbps en montée, latence ICMP $\approx$\,23--30 ms, perte $<1\%$ sur le lien 5G.
Ces outils sont directement transposables au WiFi (prochaine étape, §\ref{sec:next})
pour une comparaison côte à côte des deux réseaux.

\begin{center}
\includegraphics[width=0.9\textwidth]{img_rqt.png}
\captionof{figure}{Capture \texttt{rqt} : flux caméra Orbbec reçu sur le PC + liste des topics ROS\,2 actifs (\texttt{/cmd\_vel\_G29}, \texttt{/joy}, \texttt{/camera/color/image\_raw}, \texttt{/ticks}, \texttt{/velocity}, \texttt{/cmd\_vel\_aeb}, etc.).}
\end{center}

% =====================================================================
\section{Conduite autonome réactive : conception du module \texttt{follow\_the\_gap}}
Direction retenue : \textbf{conduite réactive} (pas de carte SLAM/Nav2 pour
l'instant). Le paquet \texttt{navigation}, jusqu'ici vide, a été créé
(\texttt{ament\_python}) avec le node \texttt{follow\_the\_gap}. Ce module a été
\emph{conçu et validé hors matériel} (cf.\ §\ref{sec:valid-synth}) pendant que le
LiDAR était indisponible, puis confirmé sur données réelles (§\ref{sec:essai}).

\subsection{Conventions de commande (rétro-ingénierie)}
À partir de \texttt{G29Teleop.py}, \texttt{TX.py} et \texttt{AEB.py} :
\begin{center}
\begin{tabularx}{\textwidth}{@{}l X@{}}
\toprule
\texttt{geometry\_msgs/Twist} & Signification (tous les topics \texttt{/cmd\_vel*}) \\
\midrule
\texttt{angular.z} & angle \textbf{servo} de direction : centre $=83$, plage $[28;138]$ ; \\
                   & $<83$ = droite, $>83$ = gauche (sens trigonométrique). \\
\texttt{linear.x}  & \textbf{vitesse} en m/s (téléop.\ avant 0,06--0,20 m/s, lent). \\
\bottomrule
\end{tabularx}
\end{center}
Le node \texttt{TX} convertit ensuite ces valeurs en PWM pour le Pico
(\texttt{/dev/ttyACM0}) via une table d'interpolation PCHIP.

\subsection{Algorithme}
\begin{enumerate}[leftmargin=1.6em]
  \item lecture de \texttt{/scan}, conservation du secteur avant ($\pm$90\textdegree),
        nettoyage (inf/NaN $\rightarrow$ portée max, lissage) ;
  \item \textbf{bulle de sécurité} autour du point le plus proche (on l'annule
        pour ne jamais viser un obstacle) ;
  \item recherche du \textbf{plus grand espace libre} (gap) ;
  \item visée du \textbf{centre du plateau le plus lointain} du gap $\rightarrow$
        cap, converti en angle servo ;
  \item \textbf{vitesse} modulée par le dégagement avant et l'intensité du
        braquage ; arrêt si l'avant $<0{,}35$ m ; watchdog d'arrêt si plus de scan
        pendant 0,5 s.
\end{enumerate}
Sortie : \texttt{/cmd\_vel\_auto}. \emph{Remarque :} viser le simple maximum
(\texttt{argmax}) faisait braquer vers le bord du gap, côté obstacle ; viser le
\emph{centre du plateau} corrige ce défaut.

\subsection{Intégration à \texttt{twist\_mux}}\label{sec:mux}
Ajout de l'entrée autonome en \textbf{priorité la plus basse}, afin que le pilote
et surtout l'AEB puissent reprendre la main à tout moment.
\begin{center}
\begin{tabular}{@{}llc@{}}
\toprule
Source & Topic & Priorité \\
\midrule
AEB (sécurité)          & \texttt{/cmd\_vel\_aeb}  & 255 \\
Pilote PS4              & \texttt{/cmd\_vel\_ps4}  & 60 \\
Pilote G29             & \texttt{/cmd\_vel\_G29}  & 30 \\
\textbf{Autonome (nouveau)} & \texttt{/cmd\_vel\_auto} & \textbf{10} \\
\bottomrule
\end{tabular}
\end{center}
Hiérarchie de sécurité : \textbf{AEB $>$ pilote $>$ autonome}.

\subsection{Validation hors matériel (scan synthétique)}\label{sec:valid-synth}
Le LiDAR étant alors indisponible, le node a été validé avec un \texttt{/scan}
simulé :
\begin{center}
\begin{tabular}{@{}lccl@{}}
\toprule
Scénario & \texttt{linear.x} & servo (\texttt{z}) & Comportement \\
\midrule
Champ libre        & 0,15 m/s & 92 ($\approx$ centre) & tout droit \\
Mur devant         & 0,0 (stop) & 138 & arrêt + braquage \\
Obstacle à droite  & 0,11 m/s & 138 (gauche) & s'écarte à gauche \\
Obstacle à gauche  & 0,10 m/s & 28 (droite)  & s'écarte à droite \\
\bottomrule
\end{tabular}
\end{center}
Tous les comportements sont conformes (\textcolor{success}{OK}), aucune erreur.

% =====================================================================
\section{Mise en service du LiDAR : diagnostic et réparation}
Cette partie est présentée comme un \textbf{journal de diagnostic}, incluant
volontairement les tentatives qui ont échoué : un LiDAR qui « tournait » mais ne
renvoyait aucune donnée. Le cheminement constitue l'essentiel du travail.

\subsection{Identification du matériel}
Inspection des périphériques de la Jetson (\texttt{lsusb}, \texttt{ls -l
/dev/ttyUSB* /dev/ttyACM*}, \texttt{/dev/serial/by-id/}) :
\begin{itemize}[leftmargin=1.4em]
  \item \texttt{ID 10c4:ea60 Silicon Labs CP210x UART Bridge} sur
        \textbf{\texttt{/dev/ttyUSB0}} : pont USB-série du LiDAR.
  \item Raspberry Pi Pico sur \texttt{/dev/ttyACM0} : carte de commande des
        actionneurs (à ne pas confondre).
  \item Caméra Orbbec Femto Bolt (USB) --- sans rapport ici.
\end{itemize}
Le pilote \texttt{rplidar\_ros} est déjà compilé dans l'espace de travail
\texttt{protova2\_yass}. Le capteur est donc un \textbf{RPLIDAR} relié par CP2102 ;
restait à déterminer le modèle, qui fixe le débit série. Détail important : le
LiDAR est monté « moteur à l'arrière », donc l'angle brut 0\textdegree{} pointe
vers l'\emph{arrière} ; on ajoute 180\textdegree{} pour ramener l'avant du robot à
0\textdegree{} (cohérent avec \texttt{detection/dist\_min.py}, et appliqué dans
\texttt{follow\_the\_gap}).

\subsection{Le symptôme : aucune donnée}
Premier essai du pilote ROS au débit 115200 :
\begin{lstlisting}
[INFO]  RPLIDAR SDK Version: 2.1.0
[ERROR] Error, operation time out. SL_RESULT_OPERATION_TIMEOUT!
[ros2run]: Process exited with failure 255
\end{lstlisting}
Le port s'ouvre, mais le capteur ne répond pas à la requête d'identification
(\texttt{getDeviceInfo}). \textbf{La tête du LiDAR tournait pourtant bien.}

\subsection{Diagnostic : élimination des causes (tentatives échouées)}
\begin{description}[leftmargin=0em, itemsep=2pt]
  \item[(\ko) Mauvais débit série.] Les quatre débits standard ont été essayés ---
        115\,200 (A1/A2), 256\,000 (A3/S1), 460\,800 (C1), 1\,000\,000 (S2/S3) ---
        tous en \emph{timeout}. Échec identique à toutes les vitesses
        $\Rightarrow$ ce n'est pas un problème de débit.
  \item[(\ko) Droits d'accès.] \texttt{id protova2 | grep dialout} $\rightarrow$
        l'utilisateur est dans le groupe \texttt{dialout} (sinon : « accès
        refusé », pas un timeout). Cause écartée.
  \item[(\ko) Interférence ModemManager.] \texttt{systemctl is-active
        ModemManager} $\rightarrow$ \texttt{inactive}. Cause écartée.
  \item[(\ko) Sonde série bas niveau.] Un script \texttt{pyserial} envoie
        directement les commandes du protocole RPLIDAR (réponse valide attendue :
        octets \texttt{A5 5A}). \textbf{Aucun octet ne revient}, à aucun débit,
        même en contournant ROS.
  \item[(\ko) Lignes DTR/RTS.] Les quatre combinaisons (avec impulsion de reset)
        donnent toutes \texttt{(rien)}. Cause écartée.
\end{description}

\subsection{Conclusion du diagnostic : panne matérielle}
\begin{table}[H]
\centering
\begin{tabularx}{\textwidth}{@{}l c X@{}}
\toprule
Élément testé & État & Interprétation \\
\midrule
Énumération USB / ouverture du port & \ok & câble USB (adaptateur$\to$Jetson) bon \\
Permissions (\texttt{dialout})       & \ok & pas un problème de droits \\
ModemManager                         & \ok & inactif, pas d'interférence \\
4 débits série                       & \ok & pas une mauvaise config \\
Lignes DTR/RTS (4 combinaisons)      & \ok & pas un souci de ligne de contrôle \\
\textbf{Réponse données du LiDAR}    & \ko & \textbf{le capteur n'émet rien} \\
Rotation de la tête (moteur)         & \ok & \textbf{alimentation présente} \\
\bottomrule
\end{tabularx}
\end{table}
\noindent \textbf{Verdict :} la tête tourne (LiDAR \emph{alimenté}) mais ne renvoie
aucune donnée $\Rightarrow$ les broches d'alimentation font contact mais
\textbf{pas les broches de données} : \textbf{défaut de câble/connecteur de
données}, ni logiciel ni alimentation. Cohérent avec le fait que ce LiDAR
fonctionnait auparavant (PrototypeVA, 2025) : connecteur probablement desserré par
les vibrations.

\begin{tcolorbox}[colback=lightgreen,colframe=success,title=Leçon clé]
La rotation de la tête ne prouve que l'\textbf{alimentation}. Une sonde série
brute muette à tous les débits prouve que le défaut est sur la \textbf{liaison de
données} (connecteur/câble), et non dans la configuration logicielle.
\end{tcolorbox}

\subsection{Réparation et validation}
Mise hors tension, puis \textbf{ré-emboîtement ferme du connecteur de données}
entre la tête du LiDAR et la carte CP2102 (tirer sur le connecteur, jamais sur les
fils ; respecter le sens), puis rebranchement de l'USB. La sonde répond alors :
\begin{lstlisting}
[115200] GET_INFO   -> a5 5a 14 00 00 00 04 18 1d 01 07 7e e4 ...
[115200] GET_HEALTH -> a5 5a 03 00 00 00 06 00 00 00
\end{lstlisting}
Décodage : descripteur \texttt{A5 5A}, \textbf{modèle 0x18}, firmware
\textbf{1.29}, révision matérielle 7, santé \texttt{00} = \textbf{OK}.
$\Rightarrow$ le capteur est un \textbf{RPLIDAR A1}, débit \textbf{115200}. Le
topic \texttt{/scan} est ensuite publié :
\begin{center}
\begin{tabular}{@{}ll@{}}
\toprule
Paramètre & Valeur \\
\midrule
Santé / mode      & OK / \emph{Sensitivity}, 8 kHz, 12 m, 10 Hz \\
Cadence \texttt{/scan} & $\approx$ 7,4 Hz \\
Points par tour   & $\approx$ 1080 \\
Portée            & 0,15 -- 12 m \\
\texttt{frame\_id} & \texttt{laser\_frame} \\
\bottomrule
\end{tabular}
\end{center}

% =====================================================================
\section{Visualisation dans RViz (côté PC, WiFi)}
RViz\,2 est fourni avec ROS\,2 Humble (\texttt{/opt/ros/humble/bin/rviz2}, déjà
installé) et une session graphique (\texttt{DISPLAY=:1}) est disponible sur le PC.

\subsection{Maintenir le LiDAR actif : tentatives échouées}
RViz tourne sur le PC ; le LiDAR doit tourner sur la Jetson. Faire survivre un
processus lancé par SSH après la fermeture de la session s'est révélé délicat :
\begin{center}
\begin{tabularx}{\textwidth}{@{}l c X@{}}
\toprule
Méthode & Résultat & Raison \\
\midrule
\texttt{nohup \ldots \& disown}      & \ko & le processus meurt à la fermeture du SSH \\
\texttt{setsid \ldots}               & \ko & idem (SIGHUP) \\
\texttt{tmux}                        & \ko & \texttt{tmux} non installé \\
SSH maintenu ouvert (corrigé)        & \ok & le LiDAR reste publié tant que la session vit \\
\bottomrule
\end{tabularx}
\end{center}
\begin{tcolorbox}[colback=lightorange,colframe=warning,title=Piège : auto-suppression par pkill]
\texttt{pkill -f "rplidar\_composition"} lancé \emph{dans} une commande SSH
correspond à la ligne de commande du shell SSH lui-même (qui contient ce texte) :
il se tue donc tout seul (sortie 255). Solutions : ménage dans une commande SSH
\emph{séparée}, astuce des crochets (\texttt{pkill -f "[r]plidar"}), ou tuer par
PID.
\end{tcolorbox}
\noindent Solution retenue : lancer le LiDAR \emph{au premier plan} dans une
session SSH maintenue ouverte ; tant qu'elle vit, \texttt{/scan} est publié.

\subsection{Configuration et résultat}
Astuce : en fixant \textbf{Fixed Frame = \texttt{laser\_frame}} (le repère du scan
lui-même), aucun arbre de transformations (TF) n'est nécessaire pour afficher le
nuage de points. Un fichier \texttt{scan.rviz} a été préparé (grille 1 m, affichage
\texttt{LaserScan} sur \texttt{/scan}, fiabilité \emph{Best Effort}), lancé par
\texttt{view\_scan\_wifi.sh}. \textbf{Résultat} (\ok) : RViz affiche en temps réel
un nuage de points rouges représentant le contour de la pièce vu de dessus ---
chaque point est un impact laser, le centre (0,0) est le LiDAR, chaque carré vaut
1 m. \texttt{LaserScan : Status Ok}, 31 fps. (« Global Status : Warn » est normal :
pas d'arbre TF, sans incidence puisque le \emph{Fixed Frame} est le repère du
scan.)

% =====================================================================
\section{Premier essai de conduite réactive (roues en l'air)}\label{sec:essai}
Une fois le \texttt{/scan} validé, un premier essai de la chaîne complète a été
réalisé. \textbf{Précaution :} la voiture était \emph{surélevée, roues ne touchant
pas le sol}, afin d'observer direction et propulsion sans déplacement réel.

\subsection{Ce qui a été lancé}
Tout est démarré par un unique script, \texttt{launch\_autonomous.sh}, qui lance
huit nodes ROS\,2 (DDS natif, domaine 125) :
\begin{table}[H]
\centering
\begin{tabularx}{\textwidth}{@{}c l X l@{}}
\toprule
\# & Node & Rôle & Topic / sortie \\
\midrule
1 & \texttt{rplidar\_composition}        & LiDAR                              & \texttt{/scan} \\
2 & \texttt{detection dist\_min}          & distances mini par secteur        & \texttt{/closest\_object\_front\_distance} \\
3 & \texttt{safety AEB}                   & freinage d'urgence (prio.\ 255)   & \texttt{/cmd\_vel\_aeb} \\
4 & \texttt{navigation follow\_the\_gap}  & conduite réactive (prio.\ 10)     & \texttt{/cmd\_vel\_auto} \\
5 & \texttt{controller g29\_teleop}       & reprise pilote (prio.\ 30)        & \texttt{/cmd\_vel\_G29} \\
6 & \texttt{twist\_mux}                   & arbitrage des priorités           & \texttt{/cmd\_vel} \\
7 & \texttt{controller TX}                & \textbf{envoi au Pico (actionneurs)} & série \texttt{/dev/ttyACM0} \\
8 & \texttt{RX}                           & retour capteurs                   & \texttt{/velocity}, \texttt{/ticks} \\
\bottomrule
\end{tabularx}
\end{table}
\noindent Hiérarchie de sécurité (arbitrée par \texttt{twist\_mux}) :
AEB (255) $>$ pilote G29 (30) $>$ autonome (10).

\subsection{Résultats observés}
La chaîne fonctionne de bout en bout. Les huit nodes sont actifs et TX a bien
ouvert le port série du Pico :
\begin{lstlisting}
[INFO] [cmd_vel_stamped_to_serial]: Serial ouvert sur /dev/ttyACM0 @ 115200
[INFO] [follow_the_gap]: cap=-82.2  servo=28.0  v=0.095 m/s  avant=1.11 m
[INFO] [cmd_vel_stamped_to_serial]: CTRL,28,0.095   <- trame envoyee au Pico
\end{lstlisting}
\begin{center}
\begin{tabular}{@{}ll@{}}
\toprule
Mesure (vue depuis le PC) & Valeur \\
\midrule
\texttt{/cmd\_vel} (servo, vitesse) & 28 ; $\approx$ 0,09 m/s \\
\texttt{/closest\_object\_front\_distance} & 0,73 m \\
follow\_the\_gap : cap visé & $\approx -82$\textdegree\ (espace libre à droite) \\
TX $\rightarrow$ Pico & \texttt{CTRL,28,0.09} en continu \\
\texttt{/velocity} (encodeur) & $-0{,}0$ \\
\bottomrule
\end{tabular}
\end{center}
\noindent \textbf{Interprétation :} le node vise le côté le plus dégagé, braque le
servo (28) et commande une vitesse lente ($\approx$0,09 m/s) --- comportement
cohérent. Les ordres atteignent physiquement le Pico. La roue motrice n'a pas
tourné car \textbf{l'ESC (moteur de propulsion) a sa propre batterie/interrupteur,
distinct de l'alimentation Jetson, et n'était pas alimenté} lors de l'essai ; la
chaîne de commande est donc correcte de bout en bout, seule la puissance de
traction manquait (d'où \texttt{/velocity}$=0$). Aucune procédure d'armement
particulière n'est requise (la conduite au G29 se fait par simple accélérateur) :
une fois l'ESC alimenté, la commande autonome entraînera la propulsion de la même
manière.

\subsection{Arrêt sécurisé}
À la demande, tout a été arrêté proprement, en garantissant l'arrêt des
actionneurs : (1) arrêt des nodes par \texttt{pkill} en \textbf{forme crochet}
(\texttt{pkill -9 -f "[r]plidar"}\ldots), sans nom littéral ailleurs dans la
commande ; (2) \textbf{envoi d'une trame neutre au Pico} pour ne pas laisser le
moteur sur la dernière commande :
\begin{lstlisting}[language=bash]
printf "CTRL,83,0.0\n" > /dev/ttyACM0   # servo=83 (centre), vitesse=0
\end{lstlisting}
\begin{tcolorbox}[colback=lightorange,colframe=warning,title=Deux pièges à l'arrêt]
\textbf{(1) Auto-suppression par \texttt{pkill} :} si la commande SSH contient un
nom de node en clair (même dans un \texttt{pgrep} de vérification), un
\texttt{pkill -f} sur ce nom tue le shell SSH lui-même (255). $\Rightarrow$ tout
en forme crochet, vérifier dans une commande séparée.\\
\textbf{(2) \texttt{kill -9} et \texttt{trap} :} tuer le script avec \texttt{-9}
n'exécute pas son \texttt{trap} de nettoyage : les nodes enfants survivent
(orphelins). $\Rightarrow$ tuer explicitement chaque node, puis envoyer la trame
neutre.
\end{tcolorbox}

% =====================================================================
\section{Livrables et scripts}
\begin{table}[H]
\centering
\begin{tabularx}{\textwidth}{@{}l X@{}}
\toprule
\textbf{Fichier} & \textbf{Description} \\
\midrule
\texttt{launch\_pc\_wifi.sh} (PC)        & démo WiFi côté PC \\
\texttt{launch\_car\_wifi.sh} (Jetson)   & démo WiFi côté véhicule \\
\texttt{view\_camera\_wifi.sh} (PC)      & visualiseur caméra (anti-plantage rqt) \\
\texttt{launch\_lidar.sh} (Jetson)       & démarrage RPLIDAR (débit paramétrable) \\
\texttt{launch\_autonomous.sh} (Jetson)  & chaîne autonome complète (8 nodes) \\
\texttt{src/navigation/} (paquet)        & \texttt{follow\_the\_gap} + launch \\
\texttt{twist\_mux.yaml}                 & ajout de l'entrée \texttt{/cmd\_vel\_auto} \\
\texttt{scan.rviz}, \texttt{view\_scan\_wifi.sh} (PC) & visualisation \texttt{/scan} \\
\texttt{rplidar\_probe.py} (Jetson)      & sonde de diagnostic série \\
\texttt{monitor\_5g.py} (Jetson)         & tableau de bord monitoring passif (port 8088) \\
\texttt{bench\_5g\_live.py} (Jetson)     & tableau de bord benchmark actif iperf3 (port 8089) \\
\bottomrule
\end{tabularx}
\end{table}
\noindent Chaîne lancée par \texttt{launch\_autonomous.sh} :
\texttt{rplidar $\rightarrow$ dist\_min $\rightarrow$ AEB \;|\; follow\_the\_gap}
$\;\rightarrow\;$ \texttt{twist\_mux} $\rightarrow$ \texttt{TX} $\rightarrow$ Pico,
avec \texttt{RX} $\rightarrow$ \texttt{/velocity}.

% =====================================================================
\section{Synthèse et prochaines étapes}\label{sec:next}
\paragraph{Synthèse.}
\begin{enumerate}[leftmargin=1.6em]
  \item Démonstration de téléopération \textbf{sur WiFi} en DDS natif validée de
        bout en bout (sans Zenoh) ; plantage \texttt{rqt} caméra corrigé.
  \item Module de \textbf{conduite réactive} \texttt{follow\_the\_gap} créé dans le
        paquet \texttt{navigation}, intégré à \texttt{twist\_mux} (priorité la plus
        basse), validé sur scan synthétique.
  \item LiDAR : \textbf{RPLIDAR A1} identifié ; panne diagnostiquée par élimination
        (\textbf{connecteur de données desserré}) ; réparé ; \texttt{/scan} publié
        ($\approx$7,4 Hz) et visualisé dans RViz.
  \item \textbf{Premier essai de conduite réactive} (roues en l'air) : chaîne
        complète opérationnelle, commandes transmises au Pico ; arrêt sécurisé.
        Reste à alimenter l'ESC pour observer la propulsion.
\end{enumerate}

\paragraph{Prochaines étapes.}
\begin{enumerate}[leftmargin=1.6em]
  \item Alimenter l'ESC (batterie de traction) et confirmer le démarrage de la
        propulsion ; ajuster si besoin la vitesse minimale (seuil de démarrage du
        moteur).
  \item Régler les paramètres de \texttt{follow\_the\_gap}
        (\texttt{bubble\_radius}, \texttt{gap\_threshold}, \texttt{steer\_gain},
        vitesses) sur données réelles.
  \item Essai au sol dans un couloir dégagé, avec reprise pilote (G29) et AEB
        actifs.
  \item À plus long terme : passage éventuel au SLAM / Nav2 (navigation
        cartographiée), qui nécessitera l'URDF (\texttt{car\_description}) et un
        arbre TF complet.
\end{enumerate}

\end{document}
